\documentclass[conference]{IEEEtran}
\IEEEoverridecommandlockouts
\usepackage{cite}
\usepackage{amsmath,amssymb}
\usepackage{graphicx}
\usepackage{booktabs}
\usepackage{xcolor}
\usepackage[hidelinks]{hyperref}

\newcommand{\NSeeds}{5}
\newcommand{\BaselinePPLShort}{123.1}
\newcommand{\BaselinePPLLong}{133.8}
\newcommand{\BaselineSDShort}{0.28}
\newcommand{\UnigramTestPPL}{942}
\newcommand{\BudgetPPLShort}{119.7}
\newcommand{\BudgetSDShort}{1.2}
\newcommand{\BudgetEpochShort}{20.3}
\newcommand{\BudgetBasePPLShort}{123.2}
\newcommand{\BudgetPPLLong}{123.0}
\newcommand{\BudgetSDLong}{1.6}
\newcommand{\BudgetEpochLong}{30.0}
\newcommand{\BudgetBasePPLLong}{133.8}


\begin{document}

\title{Which Transformer Components Matter, and Does It Depend on Context Length?\\
A Controlled Multi-Seed Ablation Study}

\author{\IEEEauthorblockN{Anonymous Authors}
\IEEEauthorblockA{Under review}}

\maketitle

%ABSTRACT%

\section{Introduction}

The Transformer~\cite{vaswani2017attention} combines several components: token and positional
embeddings, multi-head self-attention, a position-wise feed-forward network (FFN), residual
connections~\cite{he2016deep}, and layer normalization~\cite{ba2016layer}. Practitioners treat
each of these as essential, yet published ablations of them are rarely like-for-like. They
change several things at once, use a single seed, or compare models trained with different
budgets. Narang et~al.~\cite{narang2021transformer} found that many proposed Transformer
modifications fail to transfer across implementations. Bouthillier et~al.~\cite{bouthillier2021accounting}
and Dodge et~al.~\cite{dodge2020fine} showed that seed-to-seed variance alone can match or exceed
the effects that papers report.

We revisit the basic question with a deliberately narrow, tightly controlled design. Starting from
one decoder-only language model, we remove or replace exactly one component at a time, keep the
optimizer, schedule, data order, and number of training tokens fixed, and repeat every condition
over \NSeeds{} seeds. We also ask whether a component's importance depends on the
\emph{context length} the model is trained on. To answer this, we cross every condition with two
context lengths (128 and 512 tokens) while holding the number of tokens per update fixed, and we
test the length$\times$ablation interaction directly instead of comparing two separate significance
verdicts.

\textbf{Contributions.}
(i)~A full-factorial study: 8 architectural conditions $\times$ 2 context lengths $\times$ \NSeeds{} seeds
(80 training runs) under an identical compute budget, evaluated on held-out test data with
best-validation checkpoint selection.
(ii)~Effect sizes with confidence intervals and family-wise error control (Holm), including an
explicit interaction contrast for context length.
%CONTRIB3%
(iv)~An open, tested code base in which every ablation is verified to change exactly one model
field. This check matters in practice: an earlier version of our own pipeline silently ignored
ablation overrides (Section~\ref{sec:validity}).

\section{Related Work}

\textbf{Positional information.} Causal language models can recover positional information
without explicit encodings, because the causal mask breaks permutation
symmetry~\cite{irie2019language,haviv2022transformer}. Kazemnejad et~al.~\cite{kazemnejad2023impact}
found that models without positional encoding can even length-generalize better.
\textbf{Attention heads.} Many heads can be pruned after training with little
loss~\cite{michel2019sixteen,voita2019analyzing}. That evidence comes from removing heads post hoc
and does not tell us what happens when a model is \emph{trained} with fewer heads.
\textbf{FFN.} FFN layers have been interpreted as key--value memories~\cite{geva2021transformer}.
Without them, or without residual paths, pure attention degenerates toward rank
collapse~\cite{dong2021attention}.
\textbf{Normalization placement.} Pre-LN trains stably without careful warmup, whereas Post-LN
is prone to divergence and depends more on warmup~\cite{xiong2020layer,liu2020understanding}.
Our study measures all of these effects in one controlled setting.

\section{Experimental Setup}

\subsection{Model}
The baseline is a 6-layer decoder-only Transformer with $d_\text{model}=256$, 8 heads of width~32,
an FFN with inner width $d_\text{ff}=1024$ and GELU~\cite{hendrycks2016gelu}, Pre-LN residual
blocks, learned absolute position embeddings, and input/output embedding
tying~\cite{press2017using}. Dropout is 0.1 on embeddings, attention weights, the FFN hidden layer,
and each residual branch. All linear and embedding weights are initialized from
$\mathcal{N}(0,0.02^2)$. Attention uses PyTorch's fused scaled-dot-product kernel~\cite{dao2022flashattention}
with a causal mask. The model has 13.29M parameters at $L{=}128$, of which 8.52M are the tied
embedding matrix.

\subsection{Ablations}
Each ablation changes exactly one configuration field (Table~\ref{tab:main}). Our test suite
checks this property by composing every experiment configuration and diffing it against the baseline.
\begin{itemize}
\item \textbf{No positional enc.}: no position signal at all; order is visible only through the causal mask.
\item \textbf{Sinusoidal pos.\ enc.}: fixed sinusoids~\cite{vaswani2017attention} replace learned embeddings.
\item \textbf{Single head}: one 256-dimensional head; parameter count is unchanged.
\item \textbf{FFN width $1\times$}: $d_\text{ff}=256$, which removes 2.36M parameters (18\%).
\item \textbf{Linear FFN}: GELU is replaced by the identity while both projections are kept, so
  the parameter count is unchanged and the FFN is an (over-parameterized) linear map.
\item \textbf{No residual}: the skip term is removed from every sub-layer, i.e., $x\mapsto\mathrm{Dropout}(F(\mathrm{LN}(x)))$.
\item \textbf{Post-LN}: $x\mapsto\mathrm{LN}(x+\mathrm{Dropout}(F(x)))$, without the final LN.
\end{itemize}

\subsection{Data and training}
We use word-level WikiText-2~\cite{merity2017pointer}. The vocabulary (33{,}278 types) is fitted on
the training split only. Each split is tokenized on whitespace and cut into contiguous,
non-overlapping chunks of $L$ tokens (2.05M training tokens). We do not add an end-of-line token,
so absolute perplexities are not directly comparable to published WikiText-2 numbers. They are,
however, exactly comparable across our conditions. All runs use AdamW~\cite{loshchilov2019decoupled}
($\beta{=}(0.9,0.95)$, weight decay 0.1, peak learning rate $3{\times}10^{-4}$), 200 linear warm-up
steps, cosine decay to $0.1\times$ the peak, and gradient clipping at~1.0. Training lasts 20 epochs
(10{,}020 updates), with bf16 autocast and an fp32 loss.

\textbf{Context length.} For $L{=}512$ we set the batch size to 8 instead of 32. Every update
therefore sees 4{,}096 tokens at both lengths, and both lengths receive the same number of updates
and the same total number of training tokens. Only the context available to each prediction changes.

\textbf{Protocol.} After each epoch we compute validation perplexity and keep the best checkpoint.
That checkpoint is evaluated once on the test split, and test perplexity is the primary metric.
A seed fixes model initialization, dropout masks, and data order. We use the seeds
$\{42, 1337, 2023, 7, 31415\}$, and every cell of the grid uses the same seed set.

\subsection{Statistical analysis}
All comparisons use per-token test cross-entropy $\ell$, so a difference
$\Delta\ell$ corresponds to a perplexity ratio $e^{\Delta\ell}$. For each ablation and length we
compare the ablation with the baseline at the same length using Welch's
$t$-test~\cite{welch1947generalization}. The resulting 14 $p$-values are adjusted with Holm's
step-down procedure~\cite{holm1979simple}. To test whether context length modulates an effect, we
estimate the interaction contrast
\begin{equation}
I = (\bar\ell_{a,512}-\bar\ell_{b,512})-(\bar\ell_{a,128}-\bar\ell_{b,128}),
\end{equation}
where $a$ is the ablation and $b$ the baseline. Its standard error combines the four independent
cell variances, and its degrees of freedom follow Welch--Satterthwaite. We Holm-adjust the seven
interaction tests as a separate family. $I>0$ means the ablation hurts \emph{more} at longer
context, in relative (perplexity-ratio) terms. All confidence intervals are 95\%.

%RESULTS%

\section{Threats to Validity}\label{sec:validity}
\textbf{Scale.} Our models are small (13M parameters) and are trained on 2M tokens, so the
conclusions may not hold at scale; Narang et~al.~\cite{narang2021transformer} observe that
architecture effects can change with scale.
\textbf{Fixed hyperparameters.} Every condition shares the baseline's learning rate, warmup, and
schedule, so we measure each ablation \emph{as a drop-in change}, not at its own optimum. Post-LN
in particular is known to be sensitive to warmup~\cite{xiong2020layer}.
\textbf{Two lengths.} An interaction measured at two points shows that importance changes with
length, but it cannot describe the shape of the dependence.
\textbf{Pipeline errors.} During development, our training script silently ignored ablation
overrides and saved every run under the baseline's name. Those runs were discarded. We added a
regression test that composes every configuration and asserts that exactly the declared field
changes, and each result file records the fully resolved model configuration and a configuration
hash. We recommend this kind of test for any ablation code base.

%CONCLUSION%

\bibliographystyle{IEEEtran}
\bibliography{references}

\end{document}
